\section{Experimental Results}
\label{sec:results}

\subsection{Classification Performance}

Table~\ref{tab:classification} presents the per-class precision, recall,
and F1-score of the original (unpruned) GestureTCN on the test set.

\begin{table}[ht]
\centering
\caption{Classification report on the test set (original model).}
\label{tab:classification}
\begin{tabular}{@{}lcccc@{}}
\toprule
\textbf{Class} & \textbf{Precision} & \textbf{Recall} & \textbf{F1} & \textbf{Support} \\
\midrule
\texttt{grab}        & 1.000 & 0.667 & 0.800 & 6 \\
\texttt{release}     & 1.000 & 1.000 & 1.000 & 6 \\
\texttt{swipe\_up}   & 0.714 & 1.000 & 0.833 & 5 \\
\texttt{swipe\_down} & 1.000 & 1.000 & 1.000 & 5 \\
\texttt{noise}       & 0.857 & 0.857 & 0.857 & 7 \\
\midrule
\textbf{Overall accuracy} & \multicolumn{3}{c}{89.66\%} & 29 \\
\textbf{Macro F1} & \multicolumn{3}{c}{0.898} & \\
\textbf{Weighted F1} & \multicolumn{3}{c}{0.895} & \\
\bottomrule
\end{tabular}
\end{table}

The model achieves perfect performance on \texttt{release} and
\texttt{swipe\_down}, while \texttt{grab} shows lower recall (0.667)
with two samples misclassified---one as \texttt{swipe\_up} and one as
\texttt{noise}. The \texttt{swipe\_up} class also shows lower precision
(0.714) due to confusion with other classes. These errors are
understandable given the spatial and temporal similarity between these
gesture types. The overall accuracy of 89.66\% reflects the challenge
of distinguishing similar hand motions from limited test data.

\subsection{Model Compression Trade-offs}

Table~\ref{tab:compression} compares the original, pruned, and
quantized models. Structured pruning achieves 1.9$\times$ parameter
reduction while maintaining the same accuracy, demonstrating that
the original model is overparameterized for this task. INT8
quantization provides inference acceleration on mobile hardware,
though the file size slightly increases due to QDQ node overhead
for small models.

\begin{table}[ht]
\centering
\caption{Impact of model optimization on size and accuracy.}
\label{tab:compression}
\begin{tabular}{@{}lccc@{}}
\toprule
\textbf{Model} & \textbf{Params} & \textbf{Size} & \textbf{Accuracy} \\
\midrule
Original (FP32) & 87,077 & 0.34\,MB & 89.66\% \\
Pruned (FP32) & 45,877 & 0.18\,MB & 89.66\% \\
Original ONNX (FP32) & 87,077 & 0.09\,MB & 89.66\% \\
Pruned ONNX (FP32) & 45,877 & 0.09\,MB & 89.66\% \\
Pruned + Quantized (INT8) & 45,877 & 0.14\,MB & 89.66\% \\
\bottomrule
\end{tabular}
\end{table}

The ONNX format yields more compact serialization than PyTorch. The
INT8 quantized model is slightly larger than the FP32 ONNX models
because QDQ quantization inserts QuantizeLinear and DequantizeLinear
nodes around each operation. For small models ($<$100\,KB), this
metadata overhead can exceed the weight compression benefit (4$\times$
reduction from FP32 to INT8). Nevertheless, the quantized model is
selected for deployment due to its faster inference on mobile CPUs
and NPUs, where INT8 operations are typically 2--4$\times$ faster
than FP32.

Notably, all models achieve identical accuracy on the test set.
This is attributed to the limited test set size (29 samples), where
accuracy granularity is approximately 3.4\% (1/29). Future work
should evaluate on a larger test set to better distinguish model
performance differences.

\subsection{Real-Time Performance}

Table~\ref{tab:latency} summarizes the end-to-end latency breakdown in
the deployed system.

\begin{table}[ht]
\centering
\caption{Latency breakdown in the deployed two-stage pipeline.}
\label{tab:latency}
\begin{tabular}{@{}ll@{}}
\toprule
\textbf{Phase} & \textbf{Latency} \\
\midrule
IDLE stage (hand detection) & $\sim$1.0\,s (10 frames at 10\,FPS) \\
WAKEUP stage (gesture classification) & 0.5--1.5\,s \\
Network transfer (UDP broadcast) & $\sim$10\,ms \\
TCP download (screenshot) & 50--500\,ms \\
\midrule
\textbf{Total gesture-to-transfer} & 1.3--2.5\,s \\
\bottomrule
\end{tabular}
\end{table}

The system maintains low CPU utilization during IDLE ($\sim$5--8\%
single-core) and moderate utilization during WAKEUP
($\sim$15--30\% for $\leq$2~seconds), making it suitable for
continuous background operation on mobile devices.

\subsection{Demo}

GrabDrop has been tested across multiple device configurations:
Android phones paired with Linux desktops, two Android devices, and
two desktop machines. In all cases, device discovery, gesture
recognition, and screenshot transfer functioned correctly across
platforms. The source code and demo video are available in the
accompanying repository.\footnote{Repository link included in the
  project submission.}
